Transformers \citep{Vaswani17} are the underlying model behind the recent 
success of Large Language Models (LLMs). The past few years saw a large amount
of theoretical development explaining the expressive power of transformers
\citep{transformers_survey,barcelo2024logical,YangCA24,hahn20,perez,CC22,tale},
their trainability and length generalizability \citep{raspl,framework,CC22}, 
and the extent to which one can formally verify them \citep{SAL25}. 
Interestingly, it is known that transformers with fixed (finite) precision
\citep{YangCA24,barcelo2024logical,tale,LC25} recognize a well-known subclass of regular languages called
\emph{star-free languages}. Fixed-precision transformers are especially 
pertinent to real-world transformers, which are implemented on hardware with 
fixed (finite) precision.

Star-free languages form a rather small subclass of regular languages. More
precisely, a star-free regular expression allows the intersection and 
complementation operators instead of the Kleene star. For this reason, the 
regular language $a^*b^*$ is star-free because it can be defined as
$\overline{\bar{\emptyset}.b.a.\bar{\emptyset}}$. On the other hand, it is known
that regular languages like $(aa)^*$ are not star-free (cf. see
\citep{straubing-book}). This is in contrast to Recurrent Neural Networks (RNN), which can recognize all 
regular languages \citep{SS95,RNN-hierarchy}. Thus, expressivity as language
recognizers per se is perhaps not the most useful criterion for an LLM 
architecture.

In this paper, we propose \emph{succinctness} as an alternative angle in 
understanding the ``expressivity'' of transformers. More precisely, the
succinctness of a language $L$ with respect to a class $\mathcal{C}$ of 
language recognizers (e.g.  transformers, automata, etc.) measures the smallest 
(denotational) size of $T \in \mathcal{C}$ that recognizes $L$, i.e., how many
symbols are used to describe $T$. Succinctness has been studied in logic in
computer science (e.g. \citep{GS04,Stockmeyer-thesis}) as an alternative (and 
more computational)
measure of expressiveness, and has direct consequence in how computationally
difficult it is to analyze a certain expression. For example, Linear Temporal
Logic (LTL) \citep{pnueli-ltl} is expressively equivalent to star-free regular
languages (e.g. see \citep{libkin-book}), as well as a subclass of deterministic
finite automata called \emph{counter-free automata} \citep{counter-free-book}.
Despite this, it is known that LTL can be exponentially more succinct than
finite automata \citep{SistlaC85}.
%, whereas star-free regular expressions can be
%\emph{nonelementarily}\footnote{Means: cannot be founded by any towers of 
%exponentials of any fixed height} more succinct than finite automata
%(\cite{Stockmeyer-thesis}). 
In other words, certain concepts can be described considerably more succinctly 
by LTL formulas as by finite automata. This has various consequences, e.g.,
analyzing LTL formulas (e.g. checking whether they describe a trivial concept)
is provably computationally more difficult than analyzing finite automata
\citep{SistlaC85}. 

%In linguistics, succinctness has also been correlate

\paragraph{Contributions.} Our main result can be summarized as follows: 
\begin{quotation}
\emph{Transformers can describe concepts extremely succinctly.}
\end{quotation}
More precisely, we show that transformers are \emph{exponentially} more succinct than
LTL and RNN (so including state-of-the-art State-Space Models (SSMs), e.g., see
\citep{mamba,ssm-will}), and \emph{doubly exponentially} more succinct than finite automata. 
This means that, with the same descriptional size, transformers can encode
complex patterns that require exponentially (resp. doubly exponentially) larger 
descriptional sizes for LTL and RNN (resp. automata). 
As a by-product of this expressivity, one may surmise that analyzing
transformers must be computationally challenging. We show this to be the case.
That is, verifying simple properties about transformers (e.g. whether it 
recognizes a trivial language) is computationally difficult: $\EXPSPACE$-complete. 
That is, with standard complexity-theoretic assumptions, this cannot be done
in better than \emph{double exponential time}. 

In fact, we also show matching upper bound on the succinctness gap for LTL
(exponential) and automata (double exponential). %This is done by proving
%an $\EXPSPACE$ upper bound, 
That is, we provide a translation from fixed-precision transformers to
exponential-sized LTL formulas. % that runs in exponential time.
This significantly improves the previously shown doubly exponential translation by \cite{YangCA24}.
As a consequence, for any fixed-precision transformer, there is an LTL formula 
(resp.\ finite automaton) of (doubly) exponential size recognizing the same 
language. %This shows that the succinctness gaps are precise.

In proving our succinctness results, we show how transformers can count from
0 up to $2^{2^n}$, i.e., the so-called ``(doubly exponentially) large 
counters''. This
requires a subtle encoding of large counters using attention. We then prove that
the resulting languages using LTL and RNN (resp. finite automata) require
exponentially (resp. doubly exponentially) larger description.

What assumptions do we use in our results? We assume that transformers
and RNN are of a \emph{fixed (finite) precision}. This assumption is faithful to 
real-world implementations, which use only fixed-precision arithmetics. 
We also use \emph{Unique-Hard Attention Transformers (UHAT)}, which are known 
to be expressively the weakest class of transformers \citep{Hao22}, e.g., their 
languages are
known to be in a very low complexity class $\ACzero$, whereas other classes of
transformers (e.g. average-hard attention or softmax) can recognize languages
beyond $\ACzero$ (e.g. majority). Additionally, UHATs have also been used as
transformer models in theoretical works of transformers (cf.
\citep{YangCA24,tale,transformers_survey,Hao22,LC25,hahn20,barcelo2024logical,uhat-data}). 

\paragraph{Organization.} We recall some formal concepts (transformers,
automata, and logic) in \cref{sec:prelim}. We show in 
\cref{sec:emptiness} how transformers can encode an exponential tiling
problem. We show in \cref{sec:succinctness} how this implies succinctness
of UHATs relative to other representations. Applications of our results for
reasoning about transformers are discussed in \cref{sec:apps}. We
conclude the paper in \cref{sec:related}.

%One possible explanation
%to the success of transformers as LLM architecture --- despite 
%the seemingly lack of expressivity --- is that star-free languages can express 
%\emph{any} finite language. 
%
%Star-free languages are
%actually sufficiently powerful to describe \emph{any subset} $S$ of any 
%language $L$ containing words of length at most some given number $N$. Thus,
%star-free languages are as powerful for describing finite languages. 
